\section{Qualcomm QIDK Implementation}

\subsection{Hardware Qualification}

QIDK is a program and development target spanning premium Snapdragon Android platforms, not a single immutable SoC. Reproducibility therefore begins by recording the product identifier, board revision, SoC, CPU/GPU/Hexagon blocks, RAM, storage, Android build fingerprint and security patch, kernel, firmware, QAIRT/QNN version, thermal zones, power source, and device-management privileges~\cite{qualcomm_qidk}. The proposal will not transfer specifications from QCS6490 or RB5 to QIDK.

The qualification suite determines supported operators, tensor layouts, static or dynamic shapes, precision, backend fallback, initialization time, sustained thermal performance, memory, reboot recovery, and background-execution behavior. Advertised TOPS are not used as a proxy for application latency or efficiency.

\subsection{Software Components}

The Android application uses Kotlin or Java for lifecycle, health, configuration, and operator status; C++ through NDK/JNI for the QNN runtime; an encrypted bounded local database; and mutually authenticated HTTPS or MQTT to the gateway. WorkManager is reserved for deferrable work, not deterministic control timing. Continuous operation uses the managed-device facilities and foreground-service mechanisms supported by the qualified Android version.

The protocol gateway, not the application, hosts BACnet or Modbus authority. This reduces Android privileges and centralizes write-point enforcement. Application updates, models, safety policies, and feature schemas are separately signed and versioned. An atomic active/previous artifact arrangement supports rollback.

\subsection{Model Toolchain}

\begin{algorithm}[t]
	\caption{Training and QIDK packaging}
	\label{alg:packaging}
	\begin{algorithmic}[1]
		\State freeze chronological splits and feature schema
		\State train simple baselines and candidate models with fixed seeds
		\State calibrate uncertainty on validation data
		\State export candidate from PyTorch to supported ONNX representation
		\State verify framework and exported numerical outputs
		\State quantize using representative training-only calibration samples
		\State compile for available QNN targets and verify operator placement
		\State compare FP32/FP16/INT8 task metrics and calibration
		\State benchmark cold and steady latency, memory, temperature and energy
		\State create manifest with hashes, runtime, SoC and safety eligibility
		\State sign and deploy to shadow mode
	\end{algorithmic}
\end{algorithm}

The preferred path is PyTorch $\rightarrow$ ONNX $\rightarrow$ AIMET post-training quantization (PTQ) or quantization-aware training (QAT) $\rightarrow$ QAIRT/QNN compilation $\rightarrow$ Hexagon NPU. The GPU is evaluated for supported FP16 models and the CPU for preprocessing, control flow, optimization, unsupported operations, and deterministic fallback. QNN DLC may remain hardware-agnostic at the Qualcomm-runtime level; context binaries are treated as SoC-specific~\cite{qualcomm_aihub_compile}. Cloud AI Hub is optional and cannot receive sensitive proprietary models or calibration data without governance approval.

Representative calibration contains normal operations, peaks, low loads, season transitions, and known rare events. PTQ is attempted first. QAT is used only when PTQ violates an accuracy, calibration, or rare-event criterion~\cite{qualcomm_aihub_quant}. Mixed precision is not assumed to be supported by every Workbench or QNN release.

\subsection{Backend and Power Benchmark}

Each compatible backend and precision receives 100 warm-up executions and at least 1,000 measured batch-one executions. The report separates model initialization, cold inference, and thermal steady state. It includes median, p95, p99 and worst latency; peak resident memory; device temperature and frequency state; accelerator assignment; task error; and 24-hour stability.

Device power is measured with a calibrated external DC power analyzer between the supply and device input where hardware access permits. A synchronized marker brackets inference. Total and idle-subtracted energy are both reported:
\begin{equation}
	e_{inf}=\frac{\int_{t_0}^{t_1}[P_{active}(t)-P_{idle}(t)]dt}{N}.
\end{equation}
Screen, radio, camera, charge state, ambient temperature, and background tasks are fixed. Battery fuel-gauge readings may support diagnostics but do not replace external instrumentation. If the QIDK enclosure prevents direct supply measurement, the limitation and alternate measurement uncertainty are reported.

\begin{table}[t]
	\centering
	\caption{QIDK release gate for a model artifact.}
	\label{tab:qidk-gate}
	\begin{tabularx}{\columnwidth}{p{0.34\columnwidth}X}
		\toprule
		Gate          & Required evidence                                      \\
		\midrule
		Compatibility & SoC, QAIRT, firmware, schema and operator audit        \\
		Correctness   & Framework/export/device numerical comparison           \\
		Task quality  & Locked-test error, calibration and rare-event metrics  \\
		Timing        & Cold and sustained p50/p95/p99 latency                 \\
		Resources     & RAM, storage, initialization and thermal state         \\
		Energy        & Total and idle-subtracted energy per inference         \\
		Reliability   & 24-hour soak, reboot, WAN loss and corrupt-model tests \\
		Security      & Signature, hash, storage permissions and SBOM          \\
		Control       & Shadow evidence and safety approval                    \\
		\bottomrule
	\end{tabularx}
\end{table}

\subsection{Resilience}

The device caches at least 24 hours of approved schedules and the most recent valid external forecasts. Loss of WAN permits local inference and energy-only optimization only while cached inputs remain within freshness limits. Loss of the gateway, required sensor, model integrity, or time synchronization forces fallback. Reboot restores read-only operation until the service verifies artifacts, time, gateway identity, and current BMS state.
